\documentclass[12pt]{article}
\usepackage[brazil]{babel}    
\usepackage[utf8]{inputenc} 
\usepackage{amsfonts}
\usepackage{amsmath}
\usepackage{amssymb}
\usepackage{graphicx}
\usepackage{float}
\usepackage{hyperref}
\usepackage{listings}
\usepackage{xcolor}
\usepackage{booktabs}

\usepackage[top=2cm,left=2cm,right=2cm,bottom=2cm]{geometry}

\lstset{
    backgroundcolor=\color{white},
    basicstyle=\ttfamily\small,
    breaklines=true,
    captionpos=b,
    commentstyle=\color{gray},
    keywordstyle=\color{blue},
    stringstyle=\color{red},
}

\title{MC833 - Relatório Trabalho 01\\
Fase 1: Streaming de Vídeo Cliente-Servidor com Sockets Raw}
\author{Yvens Ian Prado Porto - RA 184031}

\date{}

\begin{document}
\maketitle

\section{Introdução}
Este relatório descreve o desenvolvimento da primeira fase do Trabalho 01 da disciplina MC833. O objetivo principal foi implementar uma aplicação de streaming de vídeo no modelo Cliente-Servidor utilizando \textit{Raw Sockets} na linguagem Python. Ao contornar a pilha de protocolos padrão do sistema operacional, foi necessário realizar a construção manual e o encapsulamento de todos os cabeçalhos das camadas de Rede e Transporte (IP e UDP), além da implementação do protocolo de aplicação RTP para a transmissão contínua do vídeo.

\section{Protocolo e Diagrama de Sequência}
A aplicação utiliza o protocolo {UDP} na camada de transporte, em conjunto com o protocolo {RTP} na camada de aplicação. 

A escolha do {UDP} justifica-se pela natureza do streaming de vídeo em tempo real, onde a entrega rápida dos dados é mais importante do que a confiabilidade absoluta. O TCP introduziria atrasos inaceitáveis (\textit{head-of-line blocking}) devido às retransmissões automáticas de pacotes perdidos. No entanto, como o UDP não garante ordem nem sincronização, o {RTP} foi escolhido para encapsular o vídeo. O cabeçalho RTP preenche essas lacunas ao prover um \textit{Sequence Number} (permitindo ao cliente reordenar pacotes ou detectar perdas) e um \textit{Timestamp} (essencial para a reprodução do vídeo na velocidade correta e cálculo de \textit{jitter}).

Abaixo, o Diagrama de Sequência ilustrando a comunicação:

\begin{center}
    \begin{verbatim}
        CLIENTE                                               SERVIDOR
        |                                                     |
        | ------- (UDP) "catalog" --------------------------> |
        |                                                     | (Lê diretório)
        | <------ (UDP) Lista de Vídeos (.ts) --------------- |
        |                                                     |
        | ------- (UDP) "stream video.ts" ------------------> |
        |                                                     | (Abre arquivo)
        | <------ (UDP) "[Servidor] Preparando..." ---------- |
        | <------ (UDP + RTP) Chunk 1 do Vídeo -------------- |
        | <------ (UDP + RTP) Chunk 2 do Vídeo -------------- |
        |                       ...                           |
        | <------ (UDP + RTP) Chunk N do Vídeo -------------- |
        | <------ (UDP) "[Servidor] EOF" -------------------- |
        |                                                     |
    \end{verbatim}
\end{center}

\section{Estrutura dos Cabeçalhos e Alocação de Bytes}
Para cada pacote de vídeo transmitido, a aplicação monta a seguinte estrutura de cabeçalhos (\textit{headers}) antes do envio na rede. O tamanho total de cabeçalhos por pacote é de {40 bytes}.

\begin{table}[H]
\centering
\begin{tabular}{@{}llp{8cm}@{}}
\toprule
{Protocolo} & {Tamanho} & {Campos Reservados} \\ \midrule
{IPv4}      & 20 bytes         & Version, IHL, TOS, Total Length, ID, Flags, Fragment Offset, TTL, Protocol (17), Header Checksum, Source IP, Dest IP. \\
{UDP}       & 8 bytes          & Source Port, Destination Port, Length, Checksum. \\
{RTP}       & 12 bytes         & Version, Padding, Extension, CSRC Count, Marker, Payload Type (33), Sequence Number, Timestamp, SSRC. \\ \bottomrule
\end{tabular}
\caption{Estrutura dos cabeçalhos construídos via struct do Python.}
\end{table}

\subsection{Bytes Reservados para Dados (Payload)}
A quantidade de bytes reservada exclusivamente para os dados (o vídeo em si) em cada pacote é de {1316 bytes}. Esse valor foi escolhido com base na arquitetura do \textit{MPEG Transport Stream} (.ts), onde cada pacote TS possui exatos 188 bytes. O valor de 1316 bytes permite encapsular de forma perfeitamente alinhada $7$ pacotes TS por \textit{datagrama} UDP ($188 \times 7 = 1316$), evitando fragmentação desnecessária no nível de aplicação.

\section{Análise de Desempenho e Taxa de Transmissão}
Para calcular a quantidade de pacotes por \textit{frame} e a taxa de transmissão exigida pela rede, adotamos a premissa de que os vídeos operam a uma taxa de {30 frames por segundo (fps)}.

\subsection{Metodologia de Cálculo}
\begin{enumerate}
    \item {Total de Frames:} $\text{Duração (s)} \times 30\text{ fps}$.
    \item {Bytes por Frame ($S_f$):} $\text{Tamanho do Arquivo (em bytes)} / \text{Total de Frames}$.
    \item {Pacotes por Frame ($P$):} $\lceil S_f / 1316\text{ bytes} \rceil$ (arredondado para o próximo inteiro).
    \item {Taxa de Transmissão ($T$):} Sabendo que o tamanho total de um pacote na rede (Dados + IP + UDP + RTP) é de {1356 bytes}, a taxa em bits por segundo (bps) é: $T = P \times 1356 \times 8 \times 30$.
\end{enumerate}

Para fins de precisão na conversão de armazenamento, utilizou-se $1\text{ MB} = 1.048.576\text{ bytes}$. A seguir, os resultados para os três vídeos testados no catálogo.

\subsection{Vídeo 1: edu.ts}
{Especificações:} Tamanho: 7.2 MB ($7.549.747\text{ bytes}$) | Duração: 14 s | Total de Frames: 420.
\begin{itemize}
    \item {Tamanho médio por frame:} $\approx 17.975\text{ bytes}$.
    \item {Quantidade de pacotes por frame:} $14\text{ pacotes}$.
    \item {Taxa de transmissão requerida (30fps):} $\approx 4.55\text{ Mbps}$ ($4.556.160\text{ bps}$).
\end{itemize}

\subsection{Vídeo 2: golden.ts}
{Especificações:} Tamanho: 165.1 MB ($173.119.897\text{ bytes}$) | Duração: 198 s | Total de Frames: 5940.
\begin{itemize}
    \item {Tamanho médio por frame:} $\approx 29.144\text{ bytes}$.
    \item {Quantidade de pacotes por frame:} $23\text{ pacotes}$.
    \item {Taxa de transmissão requerida (30fps):} $\approx 7.48\text{ Mbps}$ ($7.485.120\text{ bps}$).
\end{itemize}

\subsection{Vídeo 3: sioiosikis.ts}
{Especificações:} Tamanho: 78.3 MB ($82.103.500\text{ bytes}$) | Duração: 58 s | Total de Frames: 1740.
\begin{itemize}
    \item {Tamanho médio por frame:} $\approx 47.185\text{ bytes}$.
    \item {Quantidade de pacotes por frame:} $36\text{ pacotes}$.
    \item {Taxa de transmissão requerida (30fps):} $\approx 11.71\text{ Mbps}$ ($11.715.840\text{ bps}$).
\end{itemize}

\section{Resultados Finais}
O ambiente foi validado com êxito. A escolha pelo \textit{Transport Stream} foi motivada pela sua arquitetura, que diferentemente de arquivos \texttt{.mp4}, permite a leitura e reprodução contínua e sequencial sem a dependência de um índice central (\textit{moov atom}).

O sistema provou sua integridade tanto de forma passiva (verificando a semelhança absoluta de \textit{hashs} MD5 entre o arquivo origem no contêiner do Servidor e o recebido no contêiner do Cliente) quanto ativamente: foi possível iniciar um \textit{player} de vídeo apontado para o arquivo em construção na máquina \textit{host} e assistir à transmissão em tempo real simultânea ao seu processamento na rede. O sistema de tratamento de exceções do \textit{Servidor} atuou corretamente em requisições de vídeos inexistentes, evitando o travamento do \textit{Cliente}.

\section{Declaração de Uso de Inteligência Artificial}
Em conformidade com as diretrizes da disciplina, declaro que ferramentas de Inteligência Artificial Generativa foram utilizadas durante o desenvolvimento deste trabalho como assistentes de tutoria e \textit{pair-programming}.

\textbf{Justificativa e Detalhamento de Uso:}
\begin{itemize}
    \item \textbf{Empacotamento de Bits:} Auxílio na compreensão das máscaras de bits (\textit{bitwise operations}) e nos formatos da biblioteca \texttt{struct} do Python para a montagem exata dos cabeçalhos RTP, UDP e IPv4.
    \item \textbf{Depuração de Redes (\textit{Debugging}):} Suporte na identificação do comportamento de \textit{Hanging Client} (cliente travado aguardando resposta) e na lógica de extração dos 14 bytes iniciais do \textit{frame Ethernet} da camada de enlace ao utilizar o \textit{socket} \texttt{AF\_PACKET}.
    \item \textbf{Estruturação Matemática e \LaTeX:} Auxílio na formatação do relatório técnico e na montagem das equações para o cálculo formal da taxa de transmissão de rede (\textit{throughput}) e pacotes por \textit{frame}.
\end{itemize}

Ressalto que toda a arquitetura base, execução no ambiente Docker, testes de conceito, validação de integridade via \textit{hashes} MD5 e a reprodução real do fluxo de vídeo foram conduzidos, executados e interpretados de forma autoral na minha máquina local.

\end{document}